% block: shifted-systems-chapter-heading page:175 type:chapter
\clearpage\section{\BT{Applications to Shifted Linear Systems}{移位线性方程组的应用}}\label{chap:shifted-linear-systems}
\setcounter{footnote}{0}
% block: shifted-systems-abstract page:175 type:ordinary-paragraph
\noindent\textbf{\BT{Abstract}{摘要}}\quad
\BT{As seen in Sect.\ \ref{sec:computational-physics} (computational physics) and Sect.\ \ref{sec:machine-learning} (machine learning), one needs to solve a set of linear systems of the form}{如 Section\ \ref{sec:computational-physics} (计算物理) 和 Section\ \ref{sec:machine-learning} (机器学习) 所见, 需要求解一组下列形式的线性方程组}
\[(A+\sigma_k I)\bm x^{(k)}=\bm b,\quad k=1,2,\ldots,m,\]
\BT{where $\sigma_k$ is a scalar, which are called shifted linear systems. When using the LU decomposition in Sect.\ \ref{subsec:lu-decomposition}, $m$ times LU decompositions are required. Similarly, when using stationary iterative methods in Sect.\ \ref{sec:stationary-iterations}, solving $m$ linear systems, i.e., $(L+D+\sigma_k I)\bm z=\bm v$ for all $k$, is required. To solve these linear systems, one can apply a suitable Krylov subspace method to each linear system. On the other hand, if the initial residual vector with respect to the $i$th linear system and the initial residual vector with respect to the $j$th linear system are collinear, i.e., the angle between the two initial residual vectors is $0$ or $\pi$, the generated Krylov subspaces become the same, which is referred to as the shift-invariance property of Krylov subspaces. This means that we can share the information of only one Krylov subspace to solve all the shifted linear systems, leading to efficient computations of shifted linear systems. In this chapter, Krylov subspace methods using the shift-invariance property are derived from the Krylov subspace methods in Chap.\ \ref{chap:krylov-classification}.}{其中 $\sigma_k$ 为标量, 这些方程组称为移位线性方程组. 使用 Section\ \ref{subsec:lu-decomposition} 中的 LU 分解时, 需要进行 $m$ 次 LU 分解. 类似地, 使用 Section\ \ref{sec:stationary-iterations} 中的定常迭代方法时, 需要求解 $m$ 个线性方程组, 即对所有 $k$ 求解 $(L+D+\sigma_k I)\bm z=\bm v$. 为求解这些线性方程组, 可以对每个线性方程组应用适当的 Krylov 子空间方法. 另一方面, 若第 $i$ 个线性方程组的初始残差向量与第 $j$ 个线性方程组的初始残差向量共线, 即两个初始残差向量之间的夹角为 $0$ 或 $\pi$, 则生成的 Krylov 子空间相同, 这称为 Krylov 子空间的移位不变性. 这意味着可以共享仅一个 Krylov 子空间的信息来求解所有移位线性方程组, 从而高效计算移位线性方程组. 本章从 Chapter\ \ref{chap:krylov-classification} 中的 Krylov 子空间方法出发, 推导利用移位不变性的 Krylov 子空间方法.}
